% Source: https://artofproblemsolving.com/community/c3h2644119_geometric_probability_problem
% Stable identity: aops-topic-2644119-q01
Place the square at $[0,2]\times[0,2]$; the midpoints of the sides are $(1,0)$, $(2,1)$, $(1,2)$, $(0,1)$. Consider the four discs of radius $1$ about them.

\emph{Opposite midpoints} are $2$ apart, so their discs meet only at the centre $(1,1)$.

\emph{Adjacent midpoints}, e.g. $(1,0)$ and $(0,1)$, are $\sqrt2$ apart, and their discs meet in a lens with tips $(0,0)$ and $(1,1)$. Seen from either centre, the common chord subtends a right angle, so the lens consists of two circular segments with central angle $90^\circ$:
\[\text{lens}=2\left(\dfrac{\pi}{4}-\dfrac12\right)=\dfrac{\pi}{2}-1.\]
A point within $1$ of $(1,0)$ has $0\le x\le2$, and a point within $1$ of $(0,1)$ has $0\le y\le2$, so the lens lies in the square.

Two different lenses would share only points within $1$ of two opposite midpoints, i.e. only the point $(1,1)$. So the favourable region is made of four lenses (shaded) of total area $4\left(\dfrac\pi2-1\right)=2\pi-4$, and the probability is
\[\dfrac{2\pi-4}{4}=\boxed{\dfrac{\pi-2}{2}}\approx0.571.\]

\begin{asy}
usepackage("amsmath");
size(3.5cm);
defaultpen(fontsize(10pt));
path lens=arc((1,0),1,90,180)--arc((0,1),1,270,360)--cycle;
for(int i=0;i<4;++i) fill(rotate(90*i,(1,1))*lens,lightgray);
draw((0,0)--(2,0)--(2,2)--(0,2)--cycle,linewidth(0.8));
draw(arc((1,0),1,0,180),dashed);draw(arc((2,1),1,90,270),dashed);
draw(arc((1,2),1,180,360),dashed);draw(arc((0,1),1,-90,90),dashed);
dot((1,0));dot((2,1));dot((1,2));dot((0,1));
label("$2$",(1.5,0),S);label("$2$",(2,1.5),E);
\end{asy}
